% SGA 3，第 VII_B 篇，英文重建稿——第 1 轮。
% 范围：权威底本分篇第 86--87 页，第 5.1.4 节开头至公式 (3)；对命题 5.1 的应用接续于组件 81。
% 合并阅读本页码： 599--600.
% 再版印刷页码： 588--589.
% 本范围内原版页边定位： 555.
% 本范围内编者注： 149.
% 权威底本： Polo--Gille Exp7B-13oct24.pdf,
% SHA-256 D303690E0195654BBE7C3DAA5D89EA23E24CD709A3575217234AC19AE8836694.
% 曾参考 Jacob Reinhold 的 CC BY 4.0 Markdown 作为比对谱系。
% 原英文 Loop 1 阶段说明：生成图采用直接取自 PDF 的精确裁片。当前中文工程已采用原生图。

\subsubsection*{5.1.4}
\phantomsection
\label{sga3:VIIB:5.1.4}

设 \((A,J)\) 为 \(\mathcal C\) 的对象，令
\[
  (A_n)=\operatorname{Gr}_J(A)
\]
为相联系的 \(\mathbf{Alpg}/k\) 的对象；于是
\[
  A_n=\overline{J^n}/\overline{J^{n+1}}
  \qquad(n\geq0).
\]
显然代数
\[
  \mathscr A=\prod_{n\geq0}A_n
\]
由 \(A_0\) 和 \(A_1\) 生成。换言之，对 \(n\geq1\)，乘法映射
\[
  A_1\widehat\otimes_{A_0}A_1
  \widehat\otimes_{A_0}\cdots
  \widehat\otimes_{A_0}A_1
  \longrightarrow A_n
  \tag{1}
  \label{sga3:VIIB:5.1.4:eq1}
\]
为满射。

再假设 \((A,J)\) 为 \(\mathcal C\) 的余群，且 \(A\) 以及商
\[
  \overline{J^n}/\overline{J^{n+1}}
\]
在 \(k\) 上拓扑平坦。由 \SGARefLink{sga3:VIIB:5.1.2:corollary}{推论~5.1.2}，\(\operatorname{Gr}_J(A)\) 此时为 \(\mathbf{Alpg}/k\) 的余群。因此，由 \SGARefLink{sga3:VIIB:5.1.3}{5.1.3}，若置
\[
\begin{split}
  A'_n
  =
  \biggl\{
    x\in\overline{J^n}/\overline{J^{n+1}}
    \ \biggm|\
    \Delta(x)-x\widehat\otimes1
    \in
    \bigoplus_{i=1}^{n}
    &\bigl(
      \overline{J^{\,n-i}}/\overline{J^{\,n-i+1}}
    \bigr)
    \widehat\otimes_k
    \\
    &\bigl(
      \overline{J^i}/\overline{J^{i+1}}
    \bigr)
  \biggr\},
\end{split}
\tag{2}
\label{sga3:VIIB:5.1.4:eq2}
\]
则映射
\[
  A'_n\widehat\otimes_kA_0\longrightarrow A_n,
  \qquad
  a'_n\widehat\otimes a_0\longmapsto a'_na_0,
\]
为 \(\mathbf{Alpg}/k\) 中的同构。\footnote{%
\label{sga3:VIIB:nde:149}令 \(G\)、\(H\)、\(\widehat G_H\) 分别为对应于 \(A\)、\(A_0=A/J\)、\(\operatorname{Gr}_J(A)\) 的形式 \(k\)-群。则 \(\tau:A_0\hookrightarrow\operatorname{Gr}_J(A)\) 对应于包含映射 \(H\hookrightarrow\widehat G_H\) 的收缩映射 \(\pi:\widehat G_H\to H\)，而上述断言意味着 \(\widehat G_H\) 为由 \(H\) 作用于 \(N=\ker(\pi)\) 所得的半直积。}

\phantomsection
\label{sga3:VIIB:5.1.4:eq2-source-tensor-note}
\noindent\textit{底本说明。}\enspace
在此映射按元素给出的表达式中，法文 PDF 印作 \(a'_n\widehat\otimes a'_n\mapsto a'_na_0\)。定义域和乘积要求使用 \(a'_n\widehat\otimes a_0\mapsto a'_na_0\)；这里采用后者。

令 \(I_0\) 表示 \(A_0\) 的增广理想。由 \textup{\SGARefLink{sga3:VIIB:5.1.4:eq1}{(1)}} 及如下交换图，其中
\[
  A_1'{}^{\widehat\otimes n}
  =
  \underbrace{
    A'_1\widehat\otimes_k\cdots\widehat\otimes_kA'_1
  }_{n\text{ 个因子}},
\]
可知：
\begin{center}
  \phantomsection
  \label{sga3:VIIB:5.1.4:degree-one-generation-diagram}
  % SGA3-VII-DIAG-127
  % SGA3-VII-DIAG-127
% Authority: Exp7B-13oct24.pdf, local p. 87.
% Crop: figures/Exp7B_localp087_degree_one_generation_diagram.png
% Locator: tex/components/80_expose_VIIB_section514_graded_generation_en.tex, Loop 1 line 113.
\begin{tikzcd}[row sep=3em,column sep=2.1em]
A_1\widehat\otimes_{A_0}\cdots\widehat\otimes_{A_0}A_1
  \arrow[d,"m"']
  & A_1'{}^{\widehat\otimes n}\widehat\otimes_k A_0
      \arrow[l,"{\sim}"']
      \arrow[d,"{m'\widehat\otimes\operatorname{id}}"']
  & \bigl(A_1'{}^{\widehat\otimes n}\widehat\otimes_k I_0\bigr)
      \oplus A_1'{}^{\widehat\otimes n}
      \arrow[l,"{\sim}"']
      \arrow[d,"{(m'\widehat\otimes\operatorname{id})\oplus m'}"] \\
A_n
  & A'_n\widehat\otimes_k A_0 \arrow[l,"{\sim}"']
  & (A'_n\widehat\otimes_k I_0)\oplus A'_n
      \arrow[l,"{\sim}"']
\end{tikzcd}

\end{center}
乘法映射
\[
  m':
  A'_1\widehat\otimes_k\cdots\widehat\otimes_kA'_1
  \longrightarrow A'_n
  \tag{3}
  \label{sga3:VIIB:5.1.4:eq3}
\]
为满射。换言之，预有限 \(k\)-代数
\[
  \mathscr A'=\prod_{n\geq0}A'_n
\]
由其一次项生成。
